\documentclass{article}
\usepackage[margin=1in]{geometry}
\usepackage{problemset}
\showsolntrue

% Custom commands
\newcommand{\R}{\mathbb{R}}
\newcommand{\X}{\mathcal{X}}
\newcommand{\Y}{\mathcal{Y}}
\newcommand{\E}{\mathbb{E}}
\newcommand{\V}{\text{Var}}
\renewcommand{\P}{\mathbb{P}}
\newcommand{\x}{\mathbf{x}}
\newcommand{\w}{\mathbf{w}}

\begin{document}

\pstitle{%
    DSC 180A%
}{
    Participation 1%
}[due Sunday, October 4 at 11:59PM]

\begin{probset}
% --- Problem 1 ---
\begin{prob}
    Describe one thing you found interesting in the reading. Describe what it is in your own words and why you found it interesting.

    \begin{soln}
        I found it interesting that I had to SSH into the main DSMLP server before I could access my own pod. It was interesting how I had to run the launch script and how it had all the options to configure how it was launched.
    \end{soln}
\end{prob}

% --- Problem 2 ---
\begin{prob}
    Describe one thing that you found difficult to understand. Try to be specific about what you don't think you understand.

    \begin{soln}
        It was difficult accessing my team specific data because when I tried listing out my teams, it didn't show any. I asked in my Slack channel, but got no response. 
    \end{soln}
\end{prob}

% --- Problem 3 ---
\begin{prob}
    In the reading, identify one choice that the author made in their investigation and your best guess at why they made that choice. For example, this may be a choice in sub-population to study, or a choice in parameter for a model.

    \begin{soln}
        Since it wasn't a study, there aren't many choices to discuss. If I had to though, I would say one interesting choice was how we launch the pod through the launch script rather than configuring the pod through a `yaml` and using `kubectl` commands to spin up the pods.
    \end{soln}
\end{prob}
\end{probset}
\end{document}